\subsection{\texorpdfstring{应用：空间 \(\mathcal{L}_{F}^{p}\) （ \(1 \leqslant p \leqslant +\infty\) ）之间的关系}{应用：空间 L sub F to the p （ 1 <= p <= +infinity ）之间的关系}}

命题 4. --- 设 f 是以一个 Banach 空间 F 为值的可测函数；使 \(1 \leqslant p \leqslant +\infty\) 且 \(\mathrm{N}_{p}(\mathbf{f})\) 有限的数 p 所成之集 I，或为空，或为 \(\overline{R}\) 的一个区间。若 I 非空，则映射 \(p \mapsto \mathrm{N}_{p}(\mathbf{f})\) 在 I 上的限制连续；此外，若 f 非可忽略，则 \(\log \mathrm{N}_{p}(\mathbf{f})\) 是 1/p 在 \(\overline{I}\) 上的凸函数。

我们已经知道（Ch. I, No.~3, 命题 5），有限数 \(p \geqslant 1\)（使 \(\mathrm{N}_{p}(\mathbf{f}) < +\infty\) 者）所成之集 J 或为空，或为一个区间，且 \(\log \mathrm{N}_{p}(\mathbf{f})\) 是 1/p 在 J 上的凸函数（当 f 非可忽略时）；这当然蕴涵 \(p \mapsto \mathrm{N}_{p}(\mathbf{f})\) 在 J 上的连续性。

若 J 为空，则或者 \(I = \varnothing\)，或者 \(I = \{+\infty\}\)，命题在此情形是显然的；以下设 J 非空。若 f 可忽略，命题也是显然的；以下设 f 非可忽略。若 \(s \in J\)，则对每个有限数 p \textgreater{} s，\(|f|^{p} = |f|^{s}|f|^{p-s}\)，且平均不等式表明

\[
\mathrm{N} _ {p} (\mathbf {f}) \leqslant \left(\mathrm{N} _ {s} (\mathbf {f})\right) ^ {s / p} \left(\mathrm{N} _ {\infty} (\mathbf {f})\right) ^ {(p - s) / p}.\tag{12}
\]

令 \(p\) 趋于 \(+\infty\)，由此得出

\[
\limsup _ {p \to + \infty} N _ {p} (\mathbf {f}) \leqslant N _ {\infty} (\mathbf {f})  .\tag{13}
\]

这就证明了：若 \(+\infty \in I\)，则 \(J\) 含有任意大的数；于是 \(I\) 确是 \(\overline{\mathbf{R}}\) 的一个区间，且 \(\overline{I} = \overline{J}\) 。若我们能证明 \(p \mapsto N sub p(\mathbf{f})\) 在 \(\overline{J}\) 上连续，命题即告证明；而只需在 \(J\) 的端点处建立连续性。此外，我们还可以假设 \(J\) 不退缩为一点。设 \(r\) 与 \(s\) 是 \(J\) 的左、右端点（ \(r < s \leqslant +\infty\)）。设 \(A\) 是由 \(x \in X\)（满足 \(|\mathbf{f}(x)| \geqslant 1\) 者）所成之（可测）集；则

\[
\int | \mathbf {f} | ^ {p} d | \mu | = \int | \mathbf {f} | ^ {p} \varphi_ {\mathrm{A}} d | \mu | + \int | \mathbf {f} | ^ {p} \varphi_ {\mathbf {C A}} d | \mu |.
\]

当 \(p \in \mathbf{J}\) 趋于 \(r\) 时，\(|\mathbf{f}|^p\varphi_{\mathbf{A}}\) 递减地趋于 \(|\mathbf{f}|^r\varphi_{\mathbf{A}}\)，而 \(|\mathbf{f}|^p\varphi_{\mathbf{C}_{\mathbf{A}}}\) 递增地趋于 \(|\mathbf{f}|^r\varphi_{\mathbf{C}_{\mathbf{A}}}\)。于是 \(\int |\mathbf{f}|^p\varphi_{\mathbf{C}_{\mathbf{A}}}d|\mu|\) 趋于 \(\int^{*}|\mathbf{f}|^{r}\varphi_{\mathbf{C}_{\mathbf{A}}}d|\mu|\)（§1, No.~3, 定理 3）。另一方面，\(|\mathbf{f}|^p\varphi_{\mathbf{A}}\) 对 \(p \in \mathbf{J}\) 可积，且 \(\int |\mathbf{f}|^p\varphi_{\mathbf{A}}d|\mu|\) 趋于 \(\int |\mathbf{f}|^r\varphi_{\mathbf{A}}d|\mu|\)（§4, No.~3, 命题 4）。因此 \(\int |\mathbf{f}|^p d|\mu|\) 趋于 \(\int^{*}|\mathbf{f}|^{r}d|\mu|\)，这就证明了 \(p \mapsto \mathrm{N}_p(\mathbf{f})\) 在 \(r\) 处的连续性。

同样的推理可应用于点 \(s\)（若 \(s < +\infty\)）。最后，设 \(s = +\infty\) 。鉴于 (13)，只需证明

\[
\liminf _ {p \to + \infty} N _ {p} (\mathbf {f}) \geqslant N _ {\infty} (\mathbf {f}).
\]

而设 \(a\) 是使 \(0 < a < \mathrm{N}_{\infty}(\mathbf{f})\) 的一个数。由于按假设存在有限值 \(p\)（使 \(\mathrm{N}_p(\mathbf{f}) < +\infty\) 者），故集 \(A\)——由 \(x \in X\)（满足 \(|\mathbf{f}(x)| \geqslant a\) 者）组成，它可测且非可忽略——由不等式 \(\varphi_A \leqslant (|\mathbf{f}|/a)^p\) 是可积的；此外，由这一不等式推出 \(\mathrm{N}_p(\mathbf{f}) \geqslant a \cdot (|\mu|(\mathrm{A}))^{1/p}\)；令 \(p\) 趋于 \(+\infty\)，得 \(\liminf_{p \to +\infty} \mathrm{N}_p(\mathbf{f}) \geqslant a\)，证毕。

推论. --- 若 r, s, p 是满足

\[
1 \leqslant r <   p <   s \leqslant + \infty ,
\]

的三个数，则交 \(\mathcal{L}_{\mathrm{F}}^{r} \cap \mathcal{L}_{\mathrm{F}}^{s}\) 含于 \(\mathcal{L}_{\mathrm{F}}^{p}\) 。

注意，一般说来，在交 \(\mathcal{L}_{\mathrm{F}}^{r} \cap \mathcal{L}_{\mathrm{F}}^{s}\) 上由诸 \(\mathcal{L}_{\mathrm{F}}^{p}\)（ \(r < p < s\)）的拓扑诱导的拓扑互不相同。若对 \(\mu\) 不再作任何假设，则在 \(\mathcal{L}_{\mathrm{F}}^{r} \cap \mathcal{L}_{\mathrm{F}}^{s}\) 上由 \(\mathcal{L}_{\mathrm{F}}^{r}\) 与 \(\mathcal{L}_{\mathrm{F}}^{s}\) 的拓扑诱导的拓扑一般不可比较（换句话说，商 \(\mathrm{N}_{r}(\mathbf{f})/\mathrm{N}_{s}(\mathbf{f})\) 在 \(L_{F}^{r} \cap L_{F}^{s}\) 中可取到任意大的值与任意小的值；参见习题 8）。

当 \(\mu\) 是有界测度时，命题 4 可以加强：

命题 5. --- 设 \(\mu\) 是一个有界测度，\(\mathbf{f}\) 是一个 \(\mu\)-可测函数（以 Banach 空间 \(\mathbf{F}\) 为值的）。数 \(p\)（满足 \(1 \leqslant p \leqslant +\infty\) 且 \(\mathrm{N}_p(\mathbf{f})\) 有限者）所成之集 I，或为空，或为以 \(p = 1\) 为左端点且含该点的区间；此外，\((|\mu|(\mathrm{X}))^{-1/p} \mathrm{N}_p(\mathbf{f})\) 是 \(p\) 在 I 上的递增函数。

这是上面命题 4 与 Ch. I, No.~3 命题 4 的推论的直接推论。

推论. --- 若测度 \(\mu\) 有界，则关系 r < s 蕴涵 \(L_{F}^{s} \subset L_{F}^{r}\)；此外，s 阶平均收敛的拓扑细于 r 阶平均收敛的拓扑（在 \(L_{F}^{s}\) 上）。

可以证明，一般说来 \(s\) 阶平均收敛的拓扑严格细于 \(r\) 阶平均收敛的拓扑（习题 8）。

命题 6. --- 设 X 是一个离散空间，\(\mu\) 是通过在 X 的每一点放置一个质量 \(+1\) 而定义的 X 上的测度。若 \(\mathbf{f}\) 是 X 到 Banach 空间 F 中的一个映射，则数 \(p\)（满足 \(1 \leqslant p \leqslant +\infty\) 且 \(\mathrm{N}_p(\mathbf{f})\) 有限者）所成之集 I，或为空，或为以 \(+\infty\) 为右端点且含该点的区间；此外，\(\mathrm{N}_p(\mathbf{f})\) 是 \(p\) 在 I 上的递减函数。

因为，\(\mu^{*}(|\mathbf{f}|) = \sum_{x\in X}|\mathbf{f}(x)|\) 对每个函数 \(\mathbf{f}\) 成立（ \(\S 1\) , No.~1, 例），且 \(\mathrm{N}_{\infty}(\mathbf{f}) = \| \mathbf{f}\| = \sup_{x\in X}|\mathbf{f}(x)|\)；若存在数 \(\alpha >0\) 使 \(|\mathbf{f}(x)|\geqslant \alpha\) 对无穷多个 \(x\in X\) 成立，则 \(\mathrm{N}_p(\mathbf{f}) = +\infty\) 对每个有限 \(p\) 成立；在相反的情形，存在 \(x_0\in X\) 使 \(|\mathbf{f}(x_0)| = \| \mathbf{f}\|\)，于是

\[
\mathrm{N} _ {\infty} (\mathbf {f}) = | \mathbf {f} (x _ {0}) | \leqslant \mathrm{N} _ {p} (\mathbf {f})
\]

对每个有限 p 成立。~由于函数 \(\log\mathrm{N}_{p}(\mathbf{f})\) 关于 1/p 是凸的且在点 \(+\infty\) 处取其最小值，它必然是 p 在 I 上的递减函数（FRV, I, §4, No.~3, 命题 5），证毕。

推论. --- 若 X 离散且测度 \(\mu\) 由在 X 的每一点放置质量 +1 而定义，则关系 r < s 蕴涵 \(L_{F}^{r} \subset L_{F}^{s}\)；此外，r 阶平均收敛的拓扑细于 s 阶平均收敛的拓扑（在 \(L_{F}^{r}\) 上）。

